\newpage
\section{Motivação}
\label{sc:motivacao}

O ponto de partida deste trabalho é uma limitação concreta das soluções de maior adoção: o Prisma, hoje o ORM dominante no ecossistema TypeScript, não oferece suporte nativo à herança de tabela única (STI), recurso relevante para modelar hierarquias de entidades. Entre as alternativas, o TypeORM é a que mais se aproxima do modelo desejado, orientado a objetos e com herança nativa, mas mostra-se limitado no tratamento de consultas. Nenhuma das duas reúne, ao mesmo tempo, a modelagem rica e uma composição de consultas adequada. É essa lacuna que motiva a proposta: um ORM que una essas duas qualidades e que, ao fazê-lo, modernize a base sobre a qual se apoia, adotando o dialeto padronizado de decorators (Stage-3) e o runtime Bun. As consequências dessa lacuna aparecem em situações corriqueiras. Ao modelar uma hierarquia como uma entidade \emph{Pessoa} especializada em \emph{Cliente} e \emph{Funcionário}, o desenvolvedor que não dispõe de herança nativa precisa fragmentar o conceito em tabelas avulsas e recompô-lo manualmente a cada consulta, perdendo a correspondência direta entre classe e tabela ou a abstração do polimorfismo. Já o acesso ingênuo a dados associados, como percorrer uma lista de pedidos e, para cada um, ler o cliente vinculado, faz o mapeador emitir uma consulta por item (o anti-padrão N+1), de modo que uma operação aparentemente trivial degrada o desempenho à medida que os dados crescem.